\section{Single-Hint Mixed 消融}

\subsection{实验定义}

这条线现在已经有 fixed / dynamic 两个版本，它们的共同点都不是“只训练 \texttt{task1}”，而是 mixed 三任务仍然一起训练，只把 hint 限制到 \texttt{task1\_sid\_sft}：

\begin{itemize}
\item fixed single-hint mixed：只对 \texttt{task1\_sid\_sft} 注入 fixed hint，其余任务继续按 zero-hint 路径走。
\item dynamic single-hint mixed：只对 \texttt{task1\_sid\_sft} 开 dynamic cascade，其余任务被强制停在 stage-0 / no-hint 路径。
\end{itemize}

\subsection{对应脚本}

{\footnotesize
\begin{itemize}
\item 画图脚本：\path{docs/research/scripts/build_instruments_report_figures.py}
\item 脚本目录：\path{hope/Qwen2_5-3B-Isntruct-qwen4B-4-256-MIMIGenRec-grec/}
\item fixed launcher：\path{Qwen2_5-3B-Isntruct-qwen4B-4-256-MIMIGenRec-grec-rl-rule-only-fixed-hint-sid-hint-only-mixed.sh}
\item dynamic launcher：\path{Qwen2_5-3B-Isntruct-qwen4B-4-256-MIMIGenRec-grec-rl-rule-only-dynamic-hint-sid-hint-only-mixed.sh}
\end{itemize}
}

\subsection{关键指标}

\begin{table}[H]
\centering
\small
\setlength{\tabcolsep}{4pt}
\begin{tabular}{p{3.8cm} p{2.2cm} c c c p{2.4cm}}
\toprule
Variant & Best ckpt & Epoch & NDCG@10 & HR@50 & 备注 \\
\midrule
fixed taskfix & checkpoint-2997 & 1.802 & 0.0931 & 0.1941 & clean fixed 参考 \\
fixed taskfix sid-only & checkpoint-2652 & 2.000 & 0.0945 & 0.1935 & clean sid-only 上界 \\
old fixed mixed-single & checkpoint-3326 & 2.000 & 0.0953 & 0.1938 & 带 bug caveat 的历史上界 \\
fixed single-hint mixed & checkpoint-2664 & 1.602 & 0.0948 & 0.1958 & 当前最强 partial-hint 候选 \\
dynamic dual-task & checkpoint-1510 & 1.003 & 0.0930 & 0.1885 & 当前 task-filtered dynamic 参考 \\
dynamic single-hint mixed & checkpoint-2331 & 1.402 & 0.0911 & 0.1823 & 当前已同步到 checkpoint-2997，但仍弱于 dynamic baseline \\
\bottomrule
\end{tabular}
\caption{single-hint mixed 现在已经从 fixed-only 读法扩展到 fixed / dynamic 两个版本。}
\end{table}

\subsection{图表}

\begin{figure}[H]
\centering
\includegraphics[width=0.96\textwidth]{single_hint_vs_fixed_family_epoch_curves.png}
\caption{single-hint mixed vs fixed family。}
\end{figure}

\begin{figure}[H]
\centering
\includegraphics[width=0.96\textwidth]{single_hint_mixed_vs_baselines_compact_curves.png}
\caption{single-hint mixed vs dynamic family；当前 compact 图里直接放入新的 dynamic single-hint mixed。}
\end{figure}

\subsection{结论}

\begin{itemize}
\item 这条线不再只是 fixed-only 实验；现在已经能把 fixed / dynamic 两个“只对 \texttt{task1} 给 hint”的版本放到同一节里读。
\item fixed single-hint mixed 仍然是当前更强、更完整的 partial-hint 候选：它的 best 点已经同时超过 corrected \texttt{fixed taskfix sid-only} 的 \texttt{NDCG@10} 和 \texttt{HR@50}。
\item dynamic single-hint mixed 现在已经补到 \texttt{checkpoint-2997}。它的 best \texttt{NDCG@10} 点在 \texttt{checkpoint-2331 / 0.0911}，对应 \texttt{HR@50=0.1823}；peak \texttt{HR@50} 则出现在 \texttt{checkpoint-1665 / 0.1874}，而尾点回落到 `0.0905 / 0.1802`。
\item 这意味着之前在 \texttt{checkpoint-666} 看起来还算有希望的 early-window readout，并没有沿着长跑继续长成一个真正的 dynamic 候选。按当前完整得多的轨迹看，它既没有接近 fixed \texttt{single-hint mixed}，也没有超过 canonical \texttt{dynamic gather-fix} 或 \texttt{dynamic dual-task}。
\item 因而这条线现在更像“帮助理解 partial-hint 机制边界”的对照项，而不是需要优先追加资源的新主候选。真正仍值得继续解释的 partial-hint 故事，仍然主要发生在 fixed \texttt{single-hint mixed} 上。
\end{itemize}
